\chapter{Introduction}

\section{Background}

Modern web documents are fundamentally multimodal, with research papers containing critical diagrams, financial reports presenting data through visualizations, and technical documentation relying on architectural illustrations. Yet contemporary search engines remain text-centric, discarding visual content or treating images as disconnected entities. This limitation stems from historical constraints when processing images at scale was computationally prohibitive.

Recent advances have transformed what is feasible. Vision-language models like CLIP \cite{radford2021learning} and BLIP-2 \cite{li2023blip2} demonstrate effective joint reasoning over text and images through contrastive pretraining. Concurrently, large language models revolutionized document ranking through zero-shot listwise reranking, with RankGPT \cite{sun2023chatgpt} and RankZephyr \cite{pradeep2023rankzephyr} achieving state-of-the-art effectiveness by generating document permutations via prompting. However, these advances progress along parallel tracks: vision-language models focus on short image-caption matching while LLM-based rankers operate exclusively on text. No existing work performs listwise ranking over genuinely multimodal documents where relevance depends on understanding both textual and visual information. Furthermore, user interaction with search systems remains static, placing cognitive burden entirely on users for query refinement despite LLMs possessing sophisticated reasoning capabilities that could transform this paradigm \cite{lewis2020retrieval, asai2024self}.

\section{Problem Statement}

\textbf{Problem 1: Scalability Challenges for Enterprise Deployment.} Multimodal search requires expensive vision-language model inference for processing images alongside text, coordinated indexing across heterogeneous backends managing both inverted indexes and vector stores, and careful architectural optimization to achieve the subsecond query latency expected in interactive enterprise applications. Without strategic design choices balancing computational cost and ranking quality, multimodal search systems become impractical for deployment on large corporate document collections containing millions of items. The challenge lies in maintaining sophisticated cross-modal reasoning capabilities while meeting the performance requirements of production enterprise environments.

\textbf{Problem 2: Absence of Multimodal Reranking.} All existing learning-to-rank methods from classical approaches \cite{burges2005learning, cao2007learning} to recent LLM-based rerankers \cite{sun2023chatgpt, pradeep2023rankzephyr} process only text. No production-ready methods perform listwise reasoning over cross-modal evidence, assessing how document text and visual elements interact to convey information. While MMDocIR \cite{dong2025mmdocir} explored first-stage multimodal retrieval, the more impactful reranking stage remains unaddressed \cite{nogueira2019passage}. This gap prevents enterprise search systems from effectively ranking business documents where visual content quality, such as chart clarity and diagram effectiveness, significantly impacts document relevance.

\textbf{Problem 3: Limited Enterprise Knowledge Access and AI Integration.} Organizations require vertical search solutions that empower employees to efficiently access internal knowledge bases while leveraging AI assistance for comprehension and synthesis. However, existing systems lack bidirectional integration where search results improve AI-generated responses and AI-generated insights refine subsequent search queries. Traditional enterprise search presents static ranked lists, while current RAG implementations treat retrieval as one-shot preprocessing. Commercial solutions add AI summaries as passive overlays without feedback mechanisms, missing opportunities where conversational context and LLM reasoning about query intent could progressively improve retrieval quality within corporate knowledge bases through iterative refinement.

\section{Objective}

This project develops a production-ready multimodal search engine addressing these limitations through integrated solutions spanning system architecture, ranking algorithms, and conversational interaction. The objectives align with four core contributions.

\noindent\textbf{Objective 1: Scalable Multimodal Search Architecture.} Design and implement an enterprise-grade system supporting dual retrieval for text and images. The architecture employs distributed crawlers extracting content from corporate repositories, processes text queries through BM25 inverted indexes while simultaneously encoding queries via CLIP for HNSW-based image retrieval \cite{malkov2018efficient}. Key engineering challenges include distributed coordination across heterogeneous backends, incremental index updates, efficient query routing, and achieving subsecond latency for interactive enterprise applications.

\noindent\textbf{Objective 2: Iterative LLM-Guided Search Refinement.} Develop an interactive system with dual-interface presentation combining traditional ranked document lists with conversational AI summaries. The system implements a closed-loop search-RAG cycle where user follow-up questions trigger query rewriting incorporating conversation history, execute fresh multimodal retrieval and reranking, and generate refined summaries. The architecture balances query responsiveness with answer accuracy through strategic caching and incremental computation optimized for enterprise deployment.

\noindent\textbf{Objective 3: Multimodal Listwise Reranking.} Extend state-of-the-art listwise LLM ranking to multimodal business documents by adapting RankZephyr's sliding window and sorting methodology to jointly reason over text and visual content. The reranker employs vision-language models to assess cross-modal relationships including information consistency, contradiction detection, and complementarity evaluation. Setwise prompting is extended to multimodal scenarios where document comparisons jointly evaluate textual relevance and visual quality rather than independently fusing modality-specific scores.

\noindent\textbf{Objective 4: Comprehensive Empirical Evaluation.} Conduct extensive evaluation measuring multimodal reranking effectiveness across established benchmarks including TREC Deep Learning tracks \cite{craswell2020overview, craswell2021overview}, BEIR datasets \cite{thakur2021beir}, and MMDocIR \cite{dong2025mmdocir}. Baseline comparisons span BM25, CLIP-based retrieval, text-only listwise rerankers, and pairwise vision-language methods. Experiments systematically ablate visual information contribution, cross-modal interaction modeling, and architectural design choices while evaluating downstream impact on RAG-generated answer quality for enterprise knowledge retrieval tasks.

\section{Project Scope}

\subsection{System Scope}

The multimodal search engine encompasses the complete pipeline from document acquisition to result presentation, integrating retrieval, ranking, and generation components into a cohesive system.

\textbf{Distributed Web Crawling.} The system implements distributed crawlers that acquire web documents while extracting both textual content and associated images. Crawlers must respect robots.txt directives, implement polite crawling with configurable delays, handle dynamic content through JavaScript rendering when necessary, and coordinate across multiple worker nodes to maximize throughput while avoiding duplicate fetches. Extracted documents are parsed to identify text passages and linked images, with metadata including source URLs, timestamps, and document structure preserved for downstream processing.

\textbf{Hybrid Indexing Infrastructure.} Two parallel index structures support text and image retrieval respectively. The text index employs traditional inverted index data structures mapping terms to posting lists of documents containing those terms, supporting efficient BM25 scoring. The image index maintains CLIP embeddings for all extracted images in an HNSW graph structure enabling approximate nearest neighbor search with logarithmic query complexity. Both indexes must support incremental updates as new documents are crawled, with mechanisms ensuring consistency between text and image representations of the same document.

\textbf{Query Processing and Routing.} User queries are processed through both retrieval pathways simultaneously. Text queries are tokenized, stop words removed, and terms matched against the inverted index to retrieve candidate documents scored by BM25. The same query text is encoded by CLIP's text encoder and matched against image embeddings via HNSW traversal to retrieve visually relevant images with associated source documents. Results from both pathways are merged based on configurable fusion strategies before being passed to the reranking stage.

\textbf{Multimodal Reranking.} Candidate documents from first-stage retrieval, each potentially containing multiple images, are reranked using vision-language models that jointly assess textual and visual relevance. The reranker implements sliding window processing to handle top-k candidates exceeding model context limits, employing sorting algorithms to efficiently identify the most relevant documents without exhaustively comparing all pairs. Prompt engineering guides the vision-language model to evaluate cross-modal consistency, contradiction, and complementarity when scoring documents.

\textbf{Dual-Interface Presentation.} The user interface presents both a traditional ranked list of documents and a conversational panel displaying AI-generated summaries. The ranked list updates in real time as users refine queries through the conversational interface. When users pose follow-up questions, the system rewrites queries incorporating conversation history, triggers fresh retrieval and reranking, updates the document list, and generates new summaries grounded in the refined results. This closed-loop interaction continues until the user's information need is satisfied.

\subsection{Target Users}

The system is deployed within a VPN service provider organization to address knowledge management needs for internal technical teams and customer-facing support staff who require comprehensive access to product documentation, VPN technical knowledge, and censorship circumvention expertise.

\textbf{Internal Development Teams.} Software engineers and network architects require rapid access to technical documentation where network topology diagrams visually illustrate packet routing alongside textual specifications, and system metrics are presented through visualizations that complement log analysis. The multimodal search enables queries for configuration guides where architecture diagrams clearly show load balancer topologies, prioritizing documents where visual content substantively aids technical comprehension.

\textbf{Customer Service Representatives.} Support staff assisting users with VPN setup and troubleshooting require documents where screenshot sequences illustrate UI navigation steps and network diagrams show how obfuscation protocols bypass censorship. The conversational AI interface allows representatives to refine queries iteratively through natural language during live customer interactions, reducing the time cost of manual query reformulation while automatically incorporating conversation history.

\subsection{System Limitations}

\textbf{Computational and Scale Constraints.} Vision-language model inference requires GPU acceleration, limiting concurrent user capacity. The system implements periodic batch index updates rather than real-time streaming, and crawled corpus coverage focuses on technical documentation domains rather than comprehensive web coverage.

\textbf{Model Capability Boundaries.} Cross-modal understanding is bounded by vision-language models trained on image-caption pairs, which may struggle with relationships between lengthy passages and multiple images. The system is optimized for English-language documents, with potential effectiveness degradation on highly specialized domains. Evaluation benchmarks were not specifically designed for cross-modal ranking, introducing potential measurement artifacts.

\section{Project Schedule}

To be determined based on project timeline and milestone definitions.